L'hydroxyde de calcium \ce{Ca(OH)2}, ainsi que l'hydroxyde de potassium
\ce{KOH}, sont deux solides ioniques.
\begin{enumerate}
    \item Écrire l'équation des réactions de dissolution de ces deux solides.
    \item On prépare une solution d'hydroxyde de calcium en dissolvant
          $\qty{50}{\mg}$ de ce composé solide dans un volume de
          $\qty{100}{\mL}$ d'eau.
    
          Quelles sont les concentrations molaires des ions présents en
          solution~?
    \item On verse alors, dans les $\qty{100}{\mL}$ de la solution précédemment
          obtenue, $\qty{80}{\mg}$ d'hydroxyde de potassium solide.
    
          Calculer la concentration de tous les ions désormais présents en
          solution.
\end{enumerate}

